%==============================================================
%  Compte rendu -- Itération 1 -- FIP INF112
%  IMT Atlantique  |  Année 2024-2025
%  Prêt à compiler sur Overleaf (pdflatex)
%==============================================================
\documentclass[12pt,a4paper]{article}

% ---------- Encodage & langue ----------
\usepackage[utf8]{inputenc}
\usepackage[T1]{fontenc}
\usepackage[french]{babel}

% ---------- Mise en page ----------
\usepackage[top=2.5cm, bottom=2.5cm, left=2.5cm, right=2.5cm]{geometry}
\usepackage{parskip}          % espace entre paragraphes
\usepackage{microtype}        % justification fine
\usepackage{lmodern}

% ---------- Couleurs & boîtes ----------
\usepackage{xcolor}
\definecolor{imtblue}{RGB}{0,102,153}
\definecolor{lightgray}{RGB}{245,245,245}
\definecolor{darkgreen}{RGB}{0,128,0}
\definecolor{codegray}{RGB}{100,100,100}
\definecolor{codered}{RGB}{180,0,0}

% ---------- Listings (code Java) ----------
\usepackage{listings}
\lstdefinestyle{java}{
  language=Java,
  basicstyle=\ttfamily\small,
  keywordstyle=\color{imtblue}\bfseries,
  stringstyle=\color{codered},
  commentstyle=\color{darkgreen}\itshape,
  numberstyle=\tiny\color{codegray},
  numbers=left,
  numbersep=8pt,
  stepnumber=1,
  backgroundcolor=\color{lightgray},
  frame=single,
  framerule=0.5pt,
  rulecolor=\color{imtblue},
  breaklines=true,
  breakatwhitespace=false,
  tabsize=4,
  showstringspaces=false,
  captionpos=b,
  xleftmargin=15pt,
  xrightmargin=5pt,
}
\lstset{style=java}

% ---------- Graphiques & figures ----------
\usepackage{graphicx}
\usepackage{float}
\usepackage{caption}
\usepackage{subcaption}

% ---------- Tableaux ----------
\usepackage{booktabs}
\usepackage{tabularx}
\usepackage{array}
\usepackage{multirow}
\usepackage{colortbl}

% ---------- Liens & PDF ----------
\usepackage[
  colorlinks=true,
  linkcolor=imtblue,
  urlcolor=imtblue,
  citecolor=imtblue,
  pdftitle={Compte rendu -- Lot1.1 -- FIP INF112},
  pdfauthor={ZIANI / LAHOUASSA},
  pdfsubject={BibTel},
]{hyperref}

% ---------- En-têtes & pieds de page ----------
\usepackage{fancyhdr}
\pagestyle{fancy}
\fancyhf{}
\fancyhead[L]{\small\textcolor{imtblue}{\textbf{FIP INF112}}}
\fancyhead[R]{\small\textcolor{imtblue}{Lot1.1}}
\fancyfoot[C]{\small\thepage}
\renewcommand{\headrulewidth}{0.4pt}
\renewcommand{\headrule}{\hbox to\headwidth{\color{imtblue}\leaders\hrule height \headrulewidth\hfill}}

% ---------- Utilitaires ----------
\usepackage{enumitem}
\usepackage{tcolorbox}
\tcbuselibrary{skins,breakable}
\usepackage{pifont}
\newcommand{\cmark}{\textcolor{darkgreen}{\ding{51}}}
\newcommand{\xmark}{\textcolor{red}{\ding{55}}}

% ---------- Boîtes colorées ----------
\newtcolorbox{infobox}[1][]{
  colback=blue!5!white, colframe=imtblue,
  fonttitle=\bfseries, title=#1,
  breakable
}
\newtcolorbox{successbox}[1][]{
  colback=green!5!white, colframe=darkgreen,
  fonttitle=\bfseries, title=#1,
  breakable
}
\newtcolorbox{warningbox}[1][]{
  colback=yellow!10!white, colframe=orange!80!black,
  fonttitle=\bfseries, title=#1,
  breakable
}

%==============================================================
\begin{document}
%==============================================================

% ============================================================
%  PAGE DE TITRE
% ============================================================
\begin{titlepage}
  \centering
  \vspace*{1cm}

  % Ligne décorative
  {\color{imtblue}\rule{\linewidth}{2pt}}\\[0.4cm]

  {\Large\bfseries\color{imtblue} FIP INF112}\\[0.3cm]
  {\color{imtblue}\rule{\linewidth}{0.8pt}}\\[1.5cm]

  {\Huge\bfseries Compte rendu de l'itération 1}\\[0.8cm]
  {\LARGE\color{imtblue} Réseau social d'opinions -- \texttt{BibTel}}\\[0.5cm]
  {\Large\itshape Mohammed-Amine Ziani \& Wissam Lahouassa}\\[2cm]

  \vfill
  {\color{imtblue}\rule{\linewidth}{2pt}}
  \vspace{0.3cm}

  {\small IMT Atlantique -- Bretagne-Pays de la Loire -- École Mines-Télécom}
\end{titlepage}

% ============================================================
%  TABLE DES MATIÈRES
% ============================================================
\tableofcontents
\newpage

% ============================================================
%  1. INTRODUCTION
% ============================================================
\section{Introduction}

Ce document constitue le compte rendu de la première itération du projet \textit{SocialNetwork},
développé dans le cadre de l'unité d'enseignement \textbf{FIP INF112 -- GL \& Gestion de projet
informatique} à IMT Atlantique (année scolaire 2024-2025).

L'objectif de cette première itération est d'implémenter les fonctionnalités de gestion des
\textbf{membres} du réseau social BibTel, conformément aux spécifications de l'API imposée par le client
(interface \texttt{ISocialNetwork}).  Plus précisément, les méthodes visées sont~:

\begin{itemize}[noitemsep]
  \item \texttt{addMember()} -- ajout d'un nouveau membre,
  \item \texttt{nbMembers()} -- retourne le nombre de membres inscrits,
  \item \texttt{toString()} -- affichage de l'état courant du réseau.
\end{itemize}

La démarche adoptée suit le principe du \textbf{développement dirigé par les tests}~: les
scénarios de test ont été préparés \emph{avant} l'implémentation, et le code a ensuite été écrit
pour les satisfaire.


% ============================================================
%  3. CONCEPTION UML (PHASE 3)
% ============================================================
\section{Conception UML}

\subsection{Diagramme de classes}

La figure~\ref{fig:classes} présente le diagramme de classes conçu pour cette itération.
L'interface \texttt{ISocialNetwork} est implémentée par \texttt{SocialNetwork}, qui agrège
une liste de \texttt{Member}.  Les exceptions \texttt{BadEntryException} et
\texttt{MemberAlreadyExistsException} sont lancées par \texttt{SocialNetwork}.

\begin{figure}[H]
  \centering
  \includegraphics[width=1\linewidth]{classes_diagram.png}
  \caption{Diagramme de classes -- Itération 1}
  \label{fig:classes}
\end{figure}

\textbf{Points de conception notables~:}
\begin{itemize}
  \item La classe \texttt{Member} est déclarée avec une visibilité \emph{package-private} afin
        d'encapsuler les détails internes dans le package \texttt{opinion} et d'empêcher tout
        accès extérieur non contrôlé.
  \item La navigabilité de l'association \texttt{SocialNetwork} $\to$ \texttt{Member} est
        unidirectionnelle~: seul \texttt{SocialNetwork} connaît ses membres, pas l'inverse.
        Cela évite le couplage inutile.
  \item Les attributs de \texttt{Member} sont tous \texttt{private}~; les accesseurs
        (\texttt{getLogin()}, \texttt{getPassword()}, \texttt{getProfile()}) ont une visibilité
        \emph{package-private} pour la même raison d'encapsulation.
\end{itemize}

\subsection{Diagrammes de séquence}

Deux scénarios détaillés ont été modélisés.

\subsubsection{Scénario 1 -- Ajout d'un membre (cas nominal)}

\begin{figure}[H]
  \centering
  \includegraphics[width=0.88\linewidth]{sequences_diagram.png}
  \caption{Diagramme de séquence -- \texttt{addMember()} cas nominal et \texttt{toString()}}
  \label{fig:seq}
\end{figure}

Le diagramme illustre~:
\begin{enumerate}
  \item La validation des paramètres d'entrée par \texttt{SocialNetwork}.
  \item La boucle de détection de doublon (comparaison insensible à la casse et aux espaces).
  \item La création de l'instance \texttt{Member} et son insertion dans la liste interne.
  \item L'appel à \texttt{toString()} qui délègue aux méthodes d'inspection et itère sur les membres.
\end{enumerate}

% ============================================================
%  4. IMPLÉMENTATION (PHASE 4)
% ============================================================
\section{Choix d'implémentation}

\subsection{Classe \texttt{Member}}

\begin{lstlisting}[caption={Classe \texttt{Member.java} -- package \texttt{opinion}}]
package opinion;

/**
 * Represents a member of the SocialNetwork.
 */
class Member {

    private String login;
    private String password;
    private String profile;

    /**
     * @param login    the member's login
     * @param password the member's password
     * @param profile  the member's profile description
     */
    Member(String login, String password, String profile) {
        this.login    = login;
        this.password = password;
        this.profile  = profile;
    }

    String getLogin()    { return login;    }
    String getPassword() { return password; }
    String getProfile()  { return profile;  }
}
\end{lstlisting}

\textbf{Justifications des choix~:}

\begin{description}[leftmargin=2cm, labelwidth=1.8cm]
  \item[Visibilité \emph{package}] La classe est déclarée sans modificateur
        (\texttt{class Member}), ce qui la rend accessible uniquement dans le package
        \texttt{opinion}.  Cela respecte le principe d'encapsulation~: les classes extérieures
        (notamment les tests et l'IHM) ne manipulent les membres qu'au travers de l'interface
        \texttt{ISocialNetwork}.

  \item[Attributs \texttt{private}] Les trois champs (\texttt{login}, \texttt{password},
        \texttt{profile}) sont privés.  Les accesseurs correspondants ont une visibilité
        \emph{package-private}~: ils sont utilisables par \texttt{SocialNetwork} mais
        invisibles depuis l'extérieur du package.

  \item[Pas d'héritage] \texttt{Member} n'étend aucune classe et n'implémente aucune
        interface~: l'approche minimaliste recommandée par les bonnes pratiques du projet est
        respectée. \\[2cm]
\end{description}

\subsection{Classe \texttt{SocialNetwork}}

\begin{lstlisting}[caption={Classe \texttt{SocialNetwork.java} -- méthodes implémentées}]
package opinion;

import java.util.LinkedList;
import exceptions.BadEntryException;
import exceptions.MemberAlreadyExistsException;

public class SocialNetwork implements ISocialNetwork {

    private LinkedList<Member> members = new LinkedList<Member>();

    @Override
    public int nbMembers() {
        return members.size();
    }

    @Override
    public void addMember(String login, String password, String profile)
            throws BadEntryException, MemberAlreadyExistsException {

        // --- Validation des parametres ---
        if (login == null || login.trim().length() < 1)
            throw new BadEntryException("login invalide");
        if (password == null || password.trim().length() < 4)
            throw new BadEntryException("password invalide");
        if (profile == null)
            throw new BadEntryException("profile invalide");

        // --- Detection de doublon ---
        for (Member m : members) {
            if (login.trim().equalsIgnoreCase(m.getLogin().trim()))
                throw new MemberAlreadyExistsException();
        }

        // --- Ajout du nouveau membre ---
        members.add(new Member(login, password, profile));
    }

    @Override
    public String toString() {
        String s = "SocialNetwork : [" + nbMembers() + " membre(s), "
                 + nbFilms() + " film(s), " + nbBooks() + " livre(s)]\n";
        s += "Membres : ";
        for (Member m : members) {
            s += m.getLogin() + " ";
        }
        return s;
    }

}
\end{lstlisting}

\textbf{Justifications des choix~:}

\begin{description}[leftmargin=2cm, labelwidth=1.8cm]

  \item[Validation stricte] \hfill 
    \begin{itemize}[noitemsep, topsep=0pt, leftmargin=*]
      \item Le login est rejeté s'il est \texttt{null} ou si sa version sans espaces de bord
            (\texttt{trim()}) a une longueur inférieure à~1.
      \item Le mot de passe est rejeté s'il est \texttt{null} ou si \texttt{trim()} donne
            moins de~4 caractères. Cela filtre les mots de passe du type \texttt{"   qwd "}
            (3 caractères utiles).
      \item Le profil est rejeté uniquement s'il est \texttt{null}~; une chaîne vide est
            acceptée (la spécification ne l'interdit pas).
    \end{itemize}

  \item[Comparaison] \hfill \\
    L'expression \texttt{login.trim().equalsIgnoreCase(m.getLogin().trim())} garantit
    qu'un login comme \texttt{"anToine"} ou \texttt{" Antoine "} est bien détecté comme
    doublon de \texttt{"Antoine"}.

  \item[Ordre] \hfill \\
    Les vérifications \texttt{BadEntryException} sont effectuées
    \emph{avant} la boucle de détection de doublon. Ainsi, si un paramètre est invalide
    \emph{et} qu'un doublon existe, l'exception levée est \texttt{BadEntryException}, ce qui
    est cohérent avec la spécification.
    
\end{description}

% ============================================================
%  5. DOSSIER DE TESTS (PHASE 2 / PHASE 5)
% ============================================================
\section{Dossier de tests}

\subsection{Scénarios de test -- \texttt{InitTest}}

\begin{table}[H]
  \centering
  \caption{Scénarios \texttt{InitTest} -- Construction d'un réseau vide}
  \begin{tabularx}{\linewidth}{>{\ttfamily}l X c}
    \toprule
    \textbf{ID} & \textbf{Description} & \textbf{Résultat attendu} \\
    \midrule
    0.1 & \texttt{nbMembers()} sur un réseau nouvellement créé & 0 \\
    0.2 & \texttt{nbBooks()} sur un réseau nouvellement créé   & 0 \\
    0.3 & \texttt{nbFilms()} sur un réseau nouvellement créé   & 0 \\
    \bottomrule
  \end{tabularx}
\end{table}

\subsection{Scénarios de test -- \texttt{AddMemberTest}}

\subsubsection{Groupe 1 -- Paramètres invalides (\texttt{BadEntryException} attendue)}

\begin{table}[H]
  \centering
  \caption{Tests de rejet de paramètres invalides}
  \begin{tabularx}{\linewidth}{>{\ttfamily}l >{\ttfamily\small}X X}
    \toprule
    \textbf{ID} & \textbf{Appel} & \textbf{Exception attendue} \\
    \midrule
    1.1 & addMember(null, "qsdfgh", "")
        & \texttt{BadEntryException} (login null) \\
    1.2 & addMember(" ", "qsdfgh", "")
        & \texttt{BadEntryException} (login sans caractère non-espace) \\
    1.3 & addMember("B", null, "")
        & \texttt{BadEntryException} (password null) \\
    1.4 & addMember("B", "~~~qwd~", "")
        & \texttt{BadEntryException} (password < 4 car. utiles) \\
    1.5 & addMember("BBBB", "bbbb", null)
        & \texttt{BadEntryException} (profile null) \\
    \bottomrule
  \end{tabularx}

\end{table}

Pour chaque test de ce groupe, la méthode auxiliaire \texttt{addMemberBadEntryTest} vérifie
également que \texttt{nbMembers()} n'a \textbf{pas} changé après le rejet~: l'état du réseau
doit rester intact.

\subsubsection{Groupe 2a -- Ajout nominal (\texttt{addMemberOKTest})}

\begin{table}[H]
  \centering
  \caption{Tests d'ajout nominal de membres}
  \begin{tabularx}{\linewidth}{>{\ttfamily}l >{\ttfamily\small}X X}
    \toprule
    \textbf{ID} & \textbf{Appel} & \textbf{Vérification} \\
    \midrule
    2.1a & addMember("Paul", "paul", "lecteur impulsif")
         & \texttt{nbMembers()} passe de 0 à 1 \\
    2.1b & addMember("Antoine", "antoine", "grand amoureux...")
         & \texttt{nbMembers()} passe de 1 à 2 \\
    2.1c & addMember("Alice", "alice", "passionnée de BD")
         & \texttt{nbMembers()} passe de 2 à 3 \\
    \bottomrule
  \end{tabularx}
\end{table}

\subsubsection{Groupe 2b -- Doublon (\texttt{MemberAlreadyExistsException} attendue)}

\begin{table}[H]
  \centering
  \caption{Tests de détection de doublon}
  \begin{tabularx}{\linewidth}{>{\ttfamily}l >{\ttfamily\small}X X}
    \toprule
    \textbf{ID} & \textbf{Appel} & \textbf{Cas testé} \\
    \midrule
    2.2 & addMember("Paul", "abcdefghij", "")
        & Login du premier membre (doublon exact) \\
    2.3 & addMember("Alice", "abcdefghij", "")
        & Login du dernier membre (doublon exact) \\
    2.4 & addMember("anToine", "abcdefghij", "")
        & Même login, casse différente \\
    2.5 & addMember(" Antoine ", "abcdefghij", "")
        & Login encadré d'espaces \\
    2.6 & addMember("An"+"toi"+"ne", "abcdefghij", "")
        & Login obtenu par concaténation de littéraux \\
    \bottomrule
  \end{tabularx}
\end{table}

\subsubsection{Groupe 3 -- Intégrité du réseau}

\begin{table}[H]
  \centering
  \caption{Vérification de la non-modification des compteurs films/livres}
  \begin{tabularx}{\linewidth}{l X c}
    \toprule
    \textbf{ID} & \textbf{Description} & \textbf{Valeur attendue} \\
    \midrule
    3.1 & \texttt{nbFilms()} après tous les appels à \texttt{addMember()} & 0 \\
    3.2 & \texttt{nbBooks()} après tous les appels à \texttt{addMember()} & 0 \\
    \bottomrule
  \end{tabularx}
\end{table}

\subsection{Résultats d'exécution}

\begin{successbox}[Résultat global -- Suite de tests \texttt{SocialNetworkTest}]

\begin{verbatim}
Testing brand new Social Networks...
Here is the newly created SocialNetwork :
SocialNetwork : [0 membre(s), 0 film(s), 0 livre(s)]
Membres :
InitTest : 3 tests, 0 erreur(s)


**********************************************************************

Testing addMember()
Final state of the social network :
SocialNetwork : [3 membre(s), 0 film(s), 0 livre(s)]
Membres : Paul Antoine Alice
AddMemberTest : 15 tests, 0 erreur(s)


**********************************************************************

Global tests results :
18 tests, 0 erreur(s)
\end{verbatim}
\end{successbox}

\subsubsection{Analyse détaillée des résultats}

\begin{table}[H]
  \centering
  \caption{Récapitulatif des résultats par groupe de tests}
  \begin{tabularx}{\linewidth}{l X c c c}
    \toprule
    \textbf{Classe} & \textbf{Groupe} & \textbf{Tests} & \textbf{OK} & \textbf{KO} \\
    \midrule
    \texttt{InitTest}     & Construction réseau vide & 3  & 3  & 0 \\
    \multirow{4}{*}{\texttt{AddMemberTest}}
                          & BadEntry (param. invalides)  & 5  & 5  & 0 \\
                          & Ajout nominal                & 3  & 3  & 0 \\
                          & MemberAlreadyExists (doublon)& 5  & 5  & 0 \\
                          & Intégrité films/livres       & 2  & 2  & 0 \\
    \midrule
    \textbf{Total}        &                              & \textbf{18} & \textbf{18} & \textbf{0} \\
    \bottomrule
  \end{tabularx}
\end{table}

% ============================================================
%  6. OBJECTIFS ATTEINTS / NON ATTEINTS
% ============================================================
\section{Bilan de l'itération}

\begin{successbox}[Objectifs atteints -- 18/18 tests passés]
  \begin{itemize}[noitemsep]
    \item[\cmark] \texttt{nbMembers()} retourne correctement la taille de la liste interne.
    \item[\cmark] \texttt{addMember()} valide tous les paramètres d'entrée.
    \item[\cmark] \texttt{addMember()} détecte les doublons (insensible à la casse et aux espaces).
    \item[\cmark] \texttt{addMember()} crée et stocke correctement les nouveaux membres.
    \item[\cmark] \texttt{toString()} affiche l'état complet du réseau (membres, films, livres).
    \item[\cmark] Aucun effet de bord sur \texttt{nbFilms()} et \texttt{nbBooks()}.
  \end{itemize}
\end{successbox}

% ============================================================
%  8. CONCLUSION
% ============================================================
\section{Conclusion}

Cette première itération a permis d'implémenter et de valider la gestion des membres dans le
réseau social \texttt{SocialNetwork}.  L'ensemble des~18~tests de la suite automatisée passent
sans erreur, ce qui confirme la conformité du code produit aux spécifications de l'API cliente.

La démarche de développement dirigé par les tests s'est avérée efficace~: les scénarios de test,
préparés avant l'implémentation, ont guidé les choix de conception (notamment la validation des
paramètres et la comparaison insensible à la casse) et ont immédiatement mis en évidence les
erreurs lors du développement.

Les fonctionnalités relatives aux films, aux livres et aux avis seront développées dans les
itérations suivantes, en s'appuyant sur les bases consolidées dans cette itération.

\end{document}
